\documentclass[11pt]{article}

\usepackage[margin=1in]{geometry}
\usepackage[T1]{fontenc}
\usepackage{lmodern}
\usepackage{amsmath,amssymb,mathtools}
\usepackage{booktabs}
\usepackage{caption}
\usepackage{enumitem}
\usepackage{float}
\usepackage{graphicx}
\usepackage{longtable}
\usepackage{microtype}
\usepackage{hyperref}
\usepackage{xcolor}

\hypersetup{
  colorlinks=true,
  linkcolor=blue!55!black,
  urlcolor=blue!55!black,
  citecolor=blue!55!black
}

\setlist[enumerate,1]{label=(\alph*), leftmargin=2.2em}
\setlist[itemize]{leftmargin=1.6em}
\setlength{\emergencystretch}{3em}
\graphicspath{{figures/}}

\newcommand{\todoanswer}[1]{%
  \par\smallskip
  \noindent\colorbox{yellow!18}{%
    \parbox{\dimexpr\linewidth-2\fboxsep\relax}{\textbf{TODO.} #1}%
  }%
  \par\smallskip
}
\newcommand{\maybeincludegraphics}[2]{%
  \IfFileExists{#1}{%
    \includegraphics[width=\linewidth]{#1}%
  }{%
    \fbox{\parbox[c][0.28\textheight][c]{0.92\linewidth}{\centering Missing figure file: \texttt{\detokenize{#1}}\\#2}}%
  }%
}
\newcommand{\codepath}[1]{\texttt{#1}}
\newcommand{\testcmd}[1]{\texttt{#1}}

\title{CS336 Assignment 2 Systems: Writeup}
\author{}
\date{\today}

\begin{document}
\maketitle
\tableofcontents
\newpage

\section{Profiling and Benchmarking}

\subsection{Problem (benchmarking\_script): Benchmarking Script}

\begin{enumerate}
  \item \textbf{Script.}
  Implemented at Problem (benchmarking\_script).

  \item \textbf{Timings with warmup.}
  \begin{table}[H]
    \centering
    \caption{End-to-end benchmark timings with 5 warmup steps and 10 measured steps.}
    \begin{tabular}{lrrrrrr}
      \toprule
      Model & Context & Forward ms & Backward ms & Optimizer ms & Train step ms & Train std ms \\
      \midrule
      small & 512 & 23.775 & 52.262 & 7.996 & 84.032 & 3.571 \\
      medium & 512 & 70.674 & 147.435 & 20.039 & 238.148 & 5.553 \\
      large & 512 & 171.237 & 344.689 & 45.678 & 561.603 & 1.322 \\
      xl & 512 & 491.891 & 999.795 & 189.236 & 1680.922 & 4.200 \\
      10B & 512 & 1869.974 & OOM & OOM & OOM & OOM \\
      \bottomrule
    \end{tabular}
  \end{table}
  These timings were measured on one NVIDIA H100 NVL in FP32 with batch size 4 and vocabulary size 10,000; the backward and optimizer columns are computed as deltas from the forward+backward and full train-step runs. Variability was generally small relative to the mean, though the small and medium backward/train runs had a few-millisecond spread; the 10B forward pass fit, but forward+backward and full training OOMed on the 93 GiB GPU.

  \item \textbf{Timings without warmup.}
  \begin{table}[H]
    \centering
    \caption{End-to-end benchmark timings with no warmup.}
    \begin{tabular}{lrrrr}
      \toprule
      Model & Forward ms & Backward ms & Train step ms & Std ms \\
      \midrule
      small & 44.150 & 59.894 & 111.300 & 102.515 \\
      medium & 88.424 & 156.985 & 264.789 & 92.256 \\
      large & 186.402 & 353.638 & 587.421 & 88.090 \\
      xl & 508.707 & 1001.157 & 1692.820 & 70.337 \\
      \bottomrule
    \end{tabular}
  \end{table}
  Without warmup, the first measured iteration includes one-time overheads such as CUDA library initialization and kernel setup, which inflates the mean and creates much larger standard deviations than the warmed measurements. The effect is most dramatic for the smaller models, where the one-time overhead is large relative to steady-state compute time.
\end{enumerate}

\subsection{Problem (nsys\_profile): Nsight Systems Profiling}

For Nsight Systems profiling, we will use the medium and large model sizes. The context lengths are 128, 512, and the largest power-of-two context length that fit a full FP32 training step on one NVIDIA H100 NVL with batch size 4 and 5 warmup steps: 2048 for medium and 1024 for large. The next larger power of two for the large model, 2048, OOMed during attention.

\begin{table}[H]
  \centering
  \caption{Selected Nsight profiling configurations.}
  \begin{tabular}{lrrr}
    \toprule
    Model & Context 1 & Context 2 & Max fitting power-of-two context \\
    \midrule
    medium & 128 & 512 & 2048 \\
    large & 128 & 512 & 1024 \\
    \bottomrule
  \end{tabular}
\end{table}

\begin{enumerate}
  \item \textbf{Forward pass time.}
  \begin{table}[H]
    \centering
    \caption{Nsight forward-pass wall times.}
    \begin{tabular}{lrr}
      \toprule
      Model & Context length & Forward time ms \\
      \midrule
      medium & 128 & 54.426 \\
      medium & 512 & 69.802 \\
      medium & 2048 & 438.232 \\
      large & 128 & 81.218 \\
      large & 512 & 177.258 \\
      large & 1024 & 386.113 \\
      \bottomrule
    \end{tabular}
  \end{table}
  At context length 512, the Nsight timings are close to the Python benchmark timings from Problem (benchmarking\_script): 69.802 ms versus 70.674 ms for medium, and 177.258 ms versus 171.237 ms for large.

  \item \textbf{Dominant CUDA kernels.}
  \begin{table}[H]
    \centering
    \caption{Dominant CUDA kernels in the forward pass.}
    \begin{tabular}{lrlrl}
      \toprule
      Model & Context length & Forward dominant kernel & Launches & Same kernel dominates forward+backward? \\
      \midrule
      medium & 128 & \texttt{sm80\_xmma\_gemm...\_128x64...\_tn...\_cublas} & 48 & no \\
      medium & 512 & \texttt{sm80\_xmma\_gemm...\_128x128...\_tn...\_cublas} & 169 & yes \\
      medium & 2048 & \texttt{sm80\_xmma\_gemm...\_128x128...\_tn...\_cublas} & 193 & yes \\
      large & 128 & \texttt{sm80\_xmma\_gemm...\_64x64...\_tn...\_cublas} & 73 & no \\
      large & 512 & \texttt{sm80\_xmma\_gemm...\_128x64...\_tn...\_cublas} & 72 & no \\
      large & 1024 & \texttt{sm80\_xmma\_gemm...\_128x128...\_tn...\_cublas} & 145 & yes \\
      \bottomrule
    \end{tabular}
  \end{table}
  The dominant forward CUDA kernel was a cuBLAS FP32 SGEMM/XMMA matrix-multiply kernel in every profile, which means the forward pass is mostly spending its GPU time in the linear layers and attention matrix multiplications rather than in Python overhead or small elementwise operations. In the forward+backward profiles, the same exact kernel still dominates for the larger context settings, while the shorter-context runs are dominated by CUTLASS SGEMM kernels from backward gradient matrix multiplications; this is still matrix-multiply work, but with different shapes and launch implementations.

  \item \textbf{Non-matmul kernels.}
  Besides matrix multiplies, the forward pass spends noticeable CUDA time in elementwise add, multiply, and divide kernels; causal-mask \texttt{where} kernels; softmax-related \texttt{exp} and reduction kernels; normalization reductions; and tensor copy/concatenation kernels. These kernels are much smaller than GEMMs individually, but at long context lengths they become visible because attention masking and softmax operate over the full context-by-context attention matrix; for example, \texttt{medium} at context length 2048 had top non-matmul kernels around 31.437 ms for an elementwise binary operation, 26.477 ms for \texttt{where}, 25.359 ms for add, and about 14.5 ms for softmax-related multiply/\texttt{exp}.

  \item \textbf{Full training step profile.}
  \begin{table}[H]
    \centering
    \caption{Matmul-like CUDA kernel fraction in forward-only versus full training profiles.}
    \begin{tabular}{lrrr}
      \toprule
      Model & Context length & Forward matmul-like \% & Train matmul-like \% \\
      \midrule
      medium & 128 & 80.2 & 54.4 \\
      medium & 512 & 78.2 & 70.4 \\
      medium & 2048 & 64.2 & 61.9 \\
      large & 128 & 84.5 & 58.9 \\
      large & 512 & 83.2 & 74.9 \\
      large & 1024 & 77.7 & 73.2 \\
      \bottomrule
    \end{tabular}
  \end{table}
  Full training is still matrix-multiply-heavy, but it is less dominated by matmul than forward-only inference. The remaining time comes from backward-pass overhead and AdamW optimizer kernels, especially elementwise parameter updates, reductions, masking/softmax work, and gradient-related operations; the drop is largest at short context lengths where fixed non-matmul overheads are larger relative to the GEMM work.

  \item \textbf{Softmax vs. matmul in attention.}
  \begin{table}[H]
    \centering
    \caption{Forward-pass CUDA kernel time inside attention NVTX ranges.}
    \begin{tabular}{lrrrr}
      \toprule
      Model & Context length & QK scores matmul ms & Softmax ms & AV value matmul ms \\
      \midrule
      medium & 128 & 0.273 & 0.586 & 0.312 \\
      medium & 512 & 3.515 & 5.970 & 2.008 \\
      medium & 2048 & 51.766 & 89.733 & 31.467 \\
      large & 128 & 0.505 & 1.082 & 0.564 \\
      large & 512 & 6.985 & 11.968 & 4.517 \\
      large & 1024 & 24.932 & 43.566 & 15.551 \\
      \bottomrule
    \end{tabular}
  \end{table}
  In these profiles, the softmax range is comparable to, and often slightly larger than, the two attention matmul ranges combined. This is surprising from FLOP counts alone, since the QK and AV matmuls do more arithmetic by scaling with the head dimension, but our softmax/masking path is unfused and memory-, reduction-, and \texttt{exp}-heavy, while the matmuls use highly optimized GEMM kernels.
\end{enumerate}

\subsection{Problem (mixed\_precision\_accumulation): Mixed-Precision Accumulation}

\begin{verbatim}
tensor(10.0001)
tensor(9.9531, dtype=torch.float16)
tensor(10.0021)
tensor(10.0021)
\end{verbatim}

The first loop accumulates entirely in FP32, so it has only small rounding error, while the second loop accumulates in FP16 and pays the larger FP16 rounding cost at every addition, leading to a noticeably larger final error. The third and fourth loops are effectively the same for accumulation accuracy: the values being added start in FP16, but the running sum is FP32, so the repeated additions are performed in the wider accumulator and the remaining error mostly comes from the FP16 representation of the addends rather than from repeatedly accumulating in FP16.

\subsection{Problem (benchmarking\_mixed\_precision): Benchmarking Mixed Precision}

\begin{enumerate}
  \item \textbf{Toy model dtypes under FP16 autocast.}
  \begin{table}[H]
    \centering
    \caption{Toy model datatypes under FP16 autocast.}
    \begin{tabular}{ll}
      \toprule
      Component & Datatype \\
      \midrule
      Model parameters in autocast context & \texttt{torch.float32} \\
      Output of \texttt{fc1} & \texttt{torch.float16} \\
      Output of \texttt{LayerNorm} & \texttt{torch.float32} \\
      Predicted logits & \texttt{torch.float16} \\
      Loss & \texttt{torch.float32} \\
      Gradients & \texttt{torch.float32} \\
      \bottomrule
    \end{tabular}
  \end{table}

  \item \textbf{LayerNorm precision sensitivity.}
  LayerNorm is precision-sensitive mainly in the mean and variance reductions, plus the normalization step that subtracts the mean and divides by the square root of the variance. FP16 can lose accuracy or overflow/underflow in these operations, so autocast keeps LayerNorm in FP32 even though it casts linear-layer matmuls to FP16. BF16 reduces the overflow/underflow concern because it has FP32-like exponent range, but its mantissa is still low precision, so using FP32 for LayerNorm reductions can still improve numerical stability.

  \item \textbf{BF16 mixed precision benchmark.}
  \begin{table}[H]
    \centering
    \caption{FP32 vs. BF16 autocast forward/backward timing.}
    \begin{tabular}{lrrrr}
      \toprule
      Model & FP32 forward ms & BF16 forward ms & FP32 backward ms & BF16 backward ms \\
      \midrule
      small & 23.816 & 10.252 & 75.753 & 33.473 \\
      medium & 70.803 & 22.999 & 218.418 & 73.926 \\
      large & 171.272 & 44.860 & 515.350 & 146.413 \\
      xl & 490.574 & 97.248 & 1489.851 & 312.887 \\
      10B & 1865.825 & 299.748 & \textsc{OOM} & \textsc{OOM} \\
      \bottomrule
    \end{tabular}
  \end{table}
  These timings use batch size 4, context length 512, 5 warmup steps, and 10 measured steps; the backward column measures the forward, loss, and backward pass together. BF16 autocast speeds up both forward and backward substantially, with forward speedup increasing from about 2.3x for small to 6.2x for 10B, and backward speedup increasing from about 2.3x for small to 4.8x for xl. The 10B backward runs OOM in both FP32 and BF16 because autocast reduces activation and matmul precision but leaves model parameters and gradients in FP32.
\end{enumerate}

\subsection{Problem (memory\_profiling): Memory Profiling}

\begin{enumerate}
  \item \textbf{Memory profiler timelines.}
  \begin{figure}[H]
    \centering
    \maybeincludegraphics{figures/memory_xl_ctx128_forward.png}{Copy the forward-only memory\_viz screenshot here.}
    \caption{Active memory timeline for the xl model, context length 128, forward-only.}
  \end{figure}

  \begin{figure}[H]
    \centering
    \maybeincludegraphics{figures/memory_xl_ctx128_train.png}{Copy the full-training-step memory\_viz screenshot here.}
    \caption{Active memory timeline for the xl model, context length 128, full training step.}
  \end{figure}

  The peak is identified from the top of the active-memory curve in memory\_viz and cross-checked against \texttt{torch.cuda.max\_memory\_allocated()} from the benchmark JSON. The forward-only trace is mostly flat around the model-parameter footprint and peaks at 12.931 GiB, while the full training-step trace rises during the forward pass as activations are saved, peaks at 51.462 GiB, and then drops in stages during backward as saved tensors are freed.

  \item \textbf{Peak memory by context length.}
  \begin{table}[H]
    \centering
    \caption{Peak memory for xl model.}
    \begin{tabular}{lrr}
      \toprule
      Context length & Forward peak GiB & Training step peak GiB \\
      \midrule
      128 & 12.931 & 51.462 \\
      512 & 13.466 & 65.585 \\
      \bottomrule
    \end{tabular}
  \end{table}
  We use context length 512 instead of 2048 for the second row because the xl full training step at context length 2048 OOMs on the 93 GiB H100 with this implementation. Increasing the context length from 128 to 512 barely changes the forward-only peak, since inference does not save activations for backward, but it increases full-step peak memory from 51.462 GiB to 65.585 GiB because more activations are retained for the backward pass.

  \item \textbf{Mixed precision memory.}
  \begin{table}[H]
    \centering
    \caption{FP32 versus BF16 autocast peak memory for xl model.}
    \begin{tabular}{lrrrr}
      \toprule
      Context length & FP32 forward GiB & BF16 forward GiB & FP32 train GiB & BF16 train GiB \\
      \midrule
      128 & 12.931 & 19.168 & 51.462 & 51.452 \\
      512 & 13.466 & 19.426 & 65.585 & 62.718 \\
      \bottomrule
    \end{tabular}
  \end{table}
  BF16 autocast does not significantly reduce peak memory in these runs: the context-128 training peak is essentially unchanged, and the context-512 training peak drops by only about 2.9 GiB. Forward-only peak memory is actually higher under autocast, likely because model parameters remain FP32 while autocast may also keep lower-precision cached copies or temporaries; in full training, FP32 parameters, gradients, and optimizer state still dominate much of the footprint.

  \item \textbf{Residual stream activation size.}
  For the xl model in FP32, one residual-stream activation tensor has size $B \times T \times d_{\mathrm{model}} \times 4$ bytes, which is $4 \times 128 \times 2560 \times 4 / 1024^2 = 5$ MiB at context length 128 and $4 \times 512 \times 2560 \times 4 / 1024^2 = 20$ MiB at context length 512.

  \item \textbf{Largest active memory allocations.}
  At lower detail levels in memory\_viz, the largest visible thin bars for both context lengths are about 100 MiB. These are consistent with the xl model's large MLP weight matrices, since $d_{\mathrm{model}} \times d_{\mathrm{ff}} \times 4 = 2560 \times 10240 \times 4 = 100$ MiB; the stack traces point back to model parameter allocations rather than residual-stream activations, which are only 5 MiB at context length 128 and 20 MiB at context length 512.

  \item \textbf{Nsight memory profiling.}
  \todoanswer{Insert Nsight screenshots and summarize the five largest saved-for-backward contributors, percentages, and gradient tensor memory calculation.}
\end{enumerate}

\section{Single-GPU Memory}

\subsection{Problem (gradient\_checkpointing): Memory-Optimal Gradient Checkpointing}

\begin{enumerate}
  \item \textbf{Asymptotically optimal strategy.}
  I would use recursive checkpointing with roughly square-root fanout at each level. Split the $N$ blocks into about $\sqrt{N}$ contiguous chunks of size about $\sqrt{N}$, wrap each chunk in \codepath{checkpoint}, and inside each chunk repeat the same construction until the base case of a single Transformer block. This stores only the checkpoint-boundary activations at each recursive level during the forward pass, and during the backward pass it recomputes the currently-needed chunk and materializes the full residuals only for one small subproblem at a time. If $M(N)$ counts checkpoint boundary activations, the peak satisfies
  \[
    M(N) = O(\sqrt{N}) + M(\sqrt{N}) = O(\sqrt{N}),
  \]
  while the dominant saved residuals for an individual block are reduced to $O(1)$ simultaneously-live blocks. Since each recursive level recomputes all blocks once, the total compute is
  \[
    C(N) = O(N \log\log N)
  \]
  block evaluations, ignoring constant factors.

\begin{verbatim}
def recursive_checkpoint(blocks, x):
    n = len(blocks)
    if n == 1:
        return blocks[0](x)

    chunk_size = ceil(sqrt(n))
    for start in range(0, n, chunk_size):
        chunk = blocks[start:start + chunk_size]

        def run_chunk(y, chunk=chunk):
            return recursive_checkpoint(chunk, y)

        x = checkpoint(run_chunk, x, use_reentrant=False)
    return x
\end{verbatim}

  \item \textbf{Best one-level strategy for xl.}
  \begin{table}[H]
    \centering
    \caption{One-level checkpointing block-size sweep for xl, batch size 4, context length 2048. Measurements use FP32 on one H200 with one warmup step and one measured training step.}
    \begin{tabular}{lrr}
      \toprule
      Checkpoint block size & Peak memory GiB & Relative step time \\
      \midrule
      4 & 81.530 & 1.000 \\
      6 & 93.455 & 1.003 \\
      8 & 105.377 & 1.006 \\
      \bottomrule
    \end{tabular}
  \end{table}
  A first square-root heuristic suggests chunk sizes around $\sqrt{32}\approx 6$, so I compared block sizes 4, 6, and 8 as a local sweep around that guess. Among these three one-level checkpointing schedules, block size 4 had the lowest measured peak memory, while the step times were nearly identical. This suggests that, in this model, the residuals saved inside a chunk are large enough that using somewhat smaller chunks than the naive $\sqrt{N}$ guess is beneficial; the extra checkpoint-boundary activations are cheaper than materializing more Transformer blocks' saved tensors at once.
\end{enumerate}

\section{GPU Kernels}

\subsection{Problem (pytorch\_attention): PyTorch Attention Benchmarking}

\begin{table}[H]
  \centering
  \caption{PyTorch attention timings and memory on one H200. Batch size 8, single head, FP32, 5 warmup steps and 100 measured steps.}
  \resizebox{\textwidth}{!}{%
  \begin{tabular}{rrrrr}
    \toprule
    $d$ & Sequence length & Forward ms & Backward ms & Memory before backward GiB \\
    \midrule
    16 & 256 & 0.110 & 0.407 & 0.036 \\
    16 & 1024 & 0.199 & 0.567 & 0.127 \\
    16 & 4096 & 2.124 & 5.004 & 1.071 \\
    16 & 8192 & 8.591 & 19.730 & 4.079 \\
    16 & 16384 & 33.106 & 76.557 & 16.095 \\
    32 & 256 & 0.103 & 0.399 & 0.067 \\
    32 & 1024 & 0.213 & 0.576 & 0.129 \\
    32 & 4096 & 2.207 & 5.079 & 1.078 \\
    32 & 8192 & 8.929 & 20.024 & 4.094 \\
    32 & 16384 & 34.592 & 78.014 & 16.126 \\
    64 & 256 & 0.108 & 0.435 & 0.068 \\
    64 & 1024 & 0.236 & 0.580 & 0.133 \\
    64 & 4096 & 2.504 & 5.635 & 1.094 \\
    64 & 8192 & 10.182 & 22.405 & 4.126 \\
    64 & 16384 & 40.121 & 87.923 & 16.189 \\
    128 & 256 & 0.110 & 0.424 & 0.070 \\
    128 & 1024 & 0.277 & 0.662 & 0.141 \\
    128 & 4096 & 3.118 & 6.986 & 1.125 \\
    128 & 8192 & 12.547 & 26.905 & 4.188 \\
    128 & 16384 & 49.327 & 105.670 & 16.314 \\
    \bottomrule
  \end{tabular}
  }
\end{table}
The benchmark script is \codepath{scripts/benchmark\_pytorch\_attention.py}. None of these configurations OOMed on the H200, but the memory growth is clearly quadratic in sequence length and only weakly dependent on $d$: increasing from sequence length 8192 to 16384 increases the pre-backward memory by about $4\times$. For example, with batch size 8 and sequence length 16384, one FP32 attention matrix of shape $8\times 16384\times 16384$ is
\[
  8 \cdot 16384^2 \cdot 4 / 1024^3 = 8 \text{ GiB}.
\]
The measured pre-backward memory is about 16 GiB because vanilla attention materializes large sequence-by-sequence tensors such as the attention scores/probabilities for backward; the $Q,K,V$ and output activations are much smaller by comparison. The saved-for-backward attention memory therefore scales as $O(BT^2)$, which is the cost FlashAttention removes by recomputing attention probabilities tile-by-tile instead of storing the full attention matrix in HBM.

\subsection{Problem (torch\_compile): Torch Compile}

\begin{enumerate}
  \item \textbf{Compiled attention.}
	  \begin{table}[H]
	    \centering
	    \caption{Compiled vs. uncompiled PyTorch attention. Each timing cell is forward/backward ms. Speedup is computed on forward+backward time.}
	    \resizebox{\textwidth}{!}{%
	    \begin{tabular}{rrrrr}
	      \toprule
	      $d$ & Sequence length & Uncompiled forward/backward ms & Compiled forward/backward ms & Speedup \\
	      \midrule
	      16 & 256 & 0.110/0.407 & 0.075/0.307 & 1.35 \\
	      16 & 1024 & 0.199/0.567 & 0.105/0.363 & 1.64 \\
	      16 & 4096 & 2.124/5.004 & 0.814/2.159 & 2.40 \\
	      16 & 8192 & 8.591/19.730 & 2.953/8.160 & 2.55 \\
	      16 & 16384 & 33.106/76.557 & 12.064/32.044 & 2.49 \\
	      32 & 256 & 0.103/0.399 & 0.078/0.340 & 1.20 \\
	      32 & 1024 & 0.213/0.576 & 0.118/0.383 & 1.57 \\
	      32 & 4096 & 2.207/5.079 & 0.892/2.238 & 2.33 \\
	      32 & 8192 & 8.929/20.024 & 3.306/8.707 & 2.41 \\
	      32 & 16384 & 34.592/78.014 & 13.494/33.739 & 2.38 \\
	      64 & 256 & 0.108/0.435 & 0.073/0.316 & 1.40 \\
	      64 & 1024 & 0.236/0.580 & 0.142/0.401 & 1.50 \\
	      64 & 4096 & 2.504/5.635 & 1.189/2.786 & 2.05 \\
	      64 & 8192 & 10.182/22.405 & 4.542/11.098 & 2.08 \\
	      64 & 16384 & 40.121/87.923 & 19.012/43.649 & 2.04 \\
	      128 & 256 & 0.110/0.424 & 0.071/0.320 & 1.37 \\
	      128 & 1024 & 0.277/0.662 & 0.183/0.484 & 1.41 \\
	      128 & 4096 & 3.118/6.986 & 1.806/4.074 & 1.72 \\
	      128 & 8192 & 12.547/26.905 & 6.883/15.573 & 1.76 \\
	      128 & 16384 & 49.327/105.670 & 28.139/61.171 & 1.74 \\
	      \bottomrule
	    \end{tabular}
	    }
	  \end{table}
	  The compiled attention function is consistently faster after warmup, with the largest speedups on long sequences and smaller head dimensions where kernel fusion removes a substantial fraction of the eager overhead. I reset Dynamo before compiling each shape in the benchmark script to avoid recompile-limit fallback across the Cartesian-product sweep; compilation overhead itself is excluded by the warmup steps.

  \item \textbf{Compiled Transformer.}
  \begin{table}[H]
    \centering
    \caption{Compiled vs. vanilla Transformer benchmark on one H200. Context length 512, batch size 4, FP32, 5 warmup steps and 10 measured steps.}
    \resizebox{\textwidth}{!}{%
    \begin{tabular}{lrrrrrr}
      \toprule
      Model & Vanilla forward ms & Compiled forward ms & Forward speedup & Vanilla train ms & Compiled train ms & Train speedup \\
      \midrule
      small & 19.069 & 14.634 & 1.30 & 65.042 & 52.706 & 1.23 \\
      medium & 52.818 & 41.864 & 1.26 & 179.980 & 147.646 & 1.22 \\
      large & 119.204 & 99.822 & 1.19 & 409.492 & 351.741 & 1.16 \\
      xl & 328.316 & 302.604 & 1.08 & 1164.430 & 1083.091 & 1.08 \\
      \bottomrule
    \end{tabular}
    }
  \end{table}
  I reran these timings on the H200 rather than reusing the earlier H100 measurements. Compiling the full Transformer improves both forward-only and full training-step runtimes, but the speedup shrinks as the model size grows: small gets about a 1.3x forward speedup and 1.23x train-step speedup, while xl is only about 1.08x faster. The full-model gains are smaller than the isolated attention gains because much of the step is already dominated by large matmuls and optimizer work, where graph-level fusion has less room to help.
\end{enumerate}

\subsection{Problem (flash\_forward): FlashAttention-2 Forward Pass}

\begin{enumerate}
  \item \textbf{Pure PyTorch autograd function.}
  The pure PyTorch implementation is \codepath{FlashAttentionPytorchFunction} in \codepath{cs336\_systems/flash\_attention.py}. It follows the tiled online-softmax recurrence with $B_q=B_k=16$, accumulates the running numerator $O_i$, denominator $l_i$, and row maximum $m_i$ in FP32, then saves $L,Q,K,V,O$ for backward. The command \testcmd{uv run pytest -q tests/test\_attention.py -k 'test\_flash\_forward\_pass\_pytorch'} passed with 1 test passed.

  \item \textbf{Triton forward kernel.}
  The Triton forward path is implemented by \codepath{flash\_fwd\_kernel} and launched from \codepath{FlashAttentionTritonFunction.forward} in \codepath{cs336\_systems/flash\_attention.py}. The launch grid is $(\lceil N_q/16\rceil,\mathrm{batch\_size})$, so each program handles one batch element and one query tile, loads the query tile once, loops over key/value tiles, and stores one output tile plus its log-sum-exp tile. The tile sizes are $Q\_TILE\_SIZE=16$ and $K\_TILE\_SIZE=16$. The command \testcmd{uv run pytest -q tests/test\_attention.py -k 'test\_flash\_forward\_pass\_triton'} passed with 2 tests passed.

  \item \textbf{Causal masking flag.}
  Causal masking is implemented in the Triton score tile before the online softmax update. The kernel constructs absolute query and key offsets, keeps entries where $q \ge k$, and uses \codepath{tl.where} to replace masked scores with $-10^6$; it also masks key columns past $N_k$ for partial final tiles. The \codepath{is\_causal} flag is saved on the autograd context and passed as a \codepath{tl.constexpr} to both forward and backward kernels. The parametrized Triton forward tests passed for both \codepath{is\_causal=False} and \codepath{is\_causal=True}.
\end{enumerate}

\subsection{Problem (flash\_backward): FlashAttention-2 Backward Pass}

The backward implementations are in \codepath{cs336\_systems/flash\_attention.py}. The PyTorch reference backward recomputes the score and probability matrices from saved $Q,K,V,O,L$, computes $D=\mathrm{rowsum}(O \circ dO)$, and then forms $dV=P^\top dO$, $dS=P\circ(dP-D)$, $dQ=dS K/\sqrt d$, and $dK=dS^\top Q/\sqrt d$. I also implemented the optional Triton backward path following Algorithm 2: \codepath{flash\_bwd\_delta\_kernel} computes $D$, \codepath{flash\_bwd\_dkdv\_kernel} assigns one program per key tile and loops over query tiles to accumulate $dK,dV$, and \codepath{flash\_bwd\_dq\_kernel} assigns one program per query tile and loops over key tiles to accumulate $dQ$. This recomputes $P$ separately for the $dK,dV$ pass and the $dQ$ pass, avoiding atomics between programs.

The command \testcmd{uv run pytest -q tests/test\_attention.py -k 'test\_flash\_backward'} passed with 3 tests passed: the PyTorch backward test and the Triton backward tests for both causal and noncausal masking.

\subsection{Problem (flash\_benchmarking): FlashAttention-2 Benchmarking}

The benchmark script is \codepath{scripts/benchmark\_flash\_attention.py}. I used batch size 1, causal masking, $Q\_TILE\_SIZE=K\_TILE\_SIZE=16$, and \codepath{triton.testing.do\_bench} with 25 ms of warmup and 100 ms of measurement per timing. The command was \testcmd{uv run python scripts/benchmark\_flash\_attention.py --output results/flash\_benchmark\_results.csv}. The run below was on one H200, and all configurations completed without OOM.

{\scriptsize
\setlength{\tabcolsep}{3pt}
\begin{longtable}{rrrrr}
  \caption{BF16 FlashAttention-2 vs. PyTorch attention. Timing cells are forward/backward/forward+backward ms.}\\
  \toprule
  Seq. len. & $d$ & PyTorch ms & Triton ms & E2E speedup \\
  \midrule
  \endfirsthead
  \toprule
  Seq. len. & $d$ & PyTorch ms & Triton ms & E2E speedup \\
  \midrule
  \endhead
  128 & 16 & 0.072/0.287/0.557 & 0.008/0.115/0.199 & 2.80 \\
  128 & 32 & 0.070/0.282/0.580 & 0.009/0.117/0.202 & 2.87 \\
  128 & 64 & 0.070/0.289/0.610 & 0.009/0.126/0.205 & 2.98 \\
  128 & 128 & 0.078/0.304/0.604 & 0.011/0.119/0.200 & 3.02 \\
  256 & 16 & 0.077/0.300/0.571 & 0.010/0.114/0.191 & 2.99 \\
  256 & 32 & 0.072/0.294/0.589 & 0.012/0.111/0.196 & 3.01 \\
  256 & 64 & 0.075/0.296/0.567 & 0.013/0.117/0.207 & 2.74 \\
  256 & 128 & 0.070/0.299/0.584 & 0.016/0.117/0.197 & 2.96 \\
  512 & 16 & 0.072/0.287/0.581 & 0.016/0.119/0.211 & 2.75 \\
  512 & 32 & 0.070/0.284/0.555 & 0.019/0.112/0.207 & 2.68 \\
  512 & 64 & 0.068/0.299/0.594 & 0.021/0.115/0.198 & 3.00 \\
  512 & 128 & 0.071/0.296/0.597 & 0.026/0.116/0.203 & 2.94 \\
  1024 & 16 & 0.072/0.299/0.560 & 0.026/0.116/0.203 & 2.76 \\
  1024 & 32 & 0.071/0.296/0.560 & 0.034/0.115/0.209 & 2.68 \\
  1024 & 64 & 0.072/0.287/0.577 & 0.038/0.115/0.199 & 2.90 \\
  1024 & 128 & 0.073/0.290/0.663 & 0.046/0.117/0.204 & 3.25 \\
  2048 & 16 & 0.089/0.328/0.557 & 0.047/0.122/0.209 & 2.67 \\
  2048 & 32 & 0.089/0.297/0.582 & 0.065/0.119/0.200 & 2.91 \\
  2048 & 64 & 0.089/0.308/0.637 & 0.072/0.122/0.207 & 3.08 \\
  2048 & 128 & 0.090/0.303/0.584 & 0.086/0.142/0.224 & 2.61 \\
  4096 & 16 & 0.252/0.499/0.746 & 0.089/0.189/0.276 & 2.70 \\
  4096 & 32 & 0.251/0.494/0.736 & 0.131/0.175/0.303 & 2.43 \\
  4096 & 64 & 0.249/0.495/0.739 & 0.152/0.239/0.390 & 1.89 \\
  4096 & 128 & 0.252/0.498/0.743 & 0.199/0.372/0.569 & 1.31 \\
  8192 & 16 & 0.919/1.753/2.665 & 0.287/0.587/0.873 & 3.05 \\
  8192 & 32 & 0.918/1.747/2.656 & 0.355/0.551/0.902 & 2.94 \\
  8192 & 64 & 0.917/1.747/2.656 & 0.411/0.787/1.199 & 2.22 \\
  8192 & 128 & 0.924/1.754/2.669 & 0.567/1.329/1.888 & 1.41 \\
  16384 & 16 & 3.354/6.590/9.942 & 1.066/2.224/3.264 & 3.05 \\
  16384 & 32 & 3.361/6.591/9.942 & 1.285/2.169/3.443 & 2.89 \\
  16384 & 64 & 3.364/6.593/9.947 & 1.498/3.095/4.589 & 2.17 \\
  16384 & 128 & 3.383/6.617/9.988 & 2.372/5.324/7.680 & 1.30 \\
  32768 & 16 & 12.984/25.883/38.859 & 4.055/8.434/12.495 & 3.11 \\
  32768 & 32 & 13.004/25.888/38.896 & 4.995/8.471/13.453 & 2.89 \\
  32768 & 64 & 13.032/25.915/38.943 & 5.872/12.090/17.966 & 2.17 \\
  32768 & 128 & 13.079/25.977/39.032 & 8.502/20.653/29.085 & 1.34 \\
  65536 & 16 & 51.932/103.593/155.344 & 15.876/33.045/48.938 & 3.17 \\
  65536 & 32 & 51.881/103.588/155.459 & 19.672/33.513/53.179 & 2.92 \\
  65536 & 64 & 52.073/103.830/155.955 & 23.156/47.768/70.931 & 2.20 \\
  65536 & 128 & 52.348/104.085/156.790 & 32.599/81.685/114.249 & 1.37 \\
  \bottomrule
\end{longtable}
}

{\scriptsize
\setlength{\tabcolsep}{3pt}
\begin{longtable}{rrrrr}
  \caption{FP32 FlashAttention-2 vs. PyTorch attention. Timing cells are forward/backward/forward+backward ms.}\\
  \toprule
  Seq. len. & $d$ & PyTorch ms & Triton ms & E2E speedup \\
  \midrule
  \endfirsthead
  \toprule
  Seq. len. & $d$ & PyTorch ms & Triton ms & E2E speedup \\
  \midrule
  \endhead
  128 & 16 & 0.073/0.306/0.604 & 0.008/0.113/0.193 & 3.13 \\
  128 & 32 & 0.080/0.309/0.585 & 0.009/0.118/0.197 & 2.97 \\
  128 & 64 & 0.076/0.308/0.593 & 0.010/0.114/0.196 & 3.03 \\
  128 & 128 & 0.077/0.302/0.599 & 0.012/0.114/0.197 & 3.04 \\
  256 & 16 & 0.079/0.322/0.602 & 0.010/0.113/0.190 & 3.17 \\
  256 & 32 & 0.084/0.317/0.621 & 0.013/0.111/0.191 & 3.25 \\
  256 & 64 & 0.084/0.318/0.619 & 0.015/0.117/0.194 & 3.19 \\
  256 & 128 & 0.086/0.321/0.602 & 0.019/0.117/0.197 & 3.06 \\
  512 & 16 & 0.082/0.317/0.602 & 0.016/0.118/0.203 & 2.97 \\
  512 & 32 & 0.079/0.311/0.611 & 0.021/0.118/0.202 & 3.02 \\
  512 & 64 & 0.079/0.325/0.621 & 0.025/0.120/0.198 & 3.14 \\
  512 & 128 & 0.073/0.308/0.603 & 0.034/0.119/0.194 & 3.11 \\
  1024 & 16 & 0.080/0.307/0.604 & 0.029/0.122/0.211 & 2.86 \\
  1024 & 32 & 0.076/0.313/0.617 & 0.037/0.129/0.210 & 2.94 \\
  1024 & 64 & 0.076/0.320/0.620 & 0.043/0.128/0.208 & 2.98 \\
  1024 & 128 & 0.081/0.318/0.617 & 0.062/0.169/0.208 & 2.97 \\
  2048 & 16 & 0.115/0.313/0.610 & 0.054/0.173/0.202 & 3.02 \\
  2048 & 32 & 0.117/0.298/0.590 & 0.070/0.179/0.221 & 2.67 \\
  2048 & 64 & 0.132/0.311/0.600 & 0.083/0.236/0.287 & 2.09 \\
  2048 & 128 & 0.149/0.313/0.609 & 0.119/0.326/0.399 & 1.53 \\
  4096 & 16 & 0.382/0.793/1.168 & 0.110/0.350/0.454 & 2.57 \\
  4096 & 32 & 0.395/0.814/1.201 & 0.152/0.373/0.518 & 2.32 \\
  4096 & 64 & 0.438/0.880/1.311 & 0.194/0.542/0.734 & 1.79 \\
  4096 & 128 & 0.522/1.037/1.558 & 0.293/0.877/1.163 & 1.34 \\
  8192 & 16 & 1.385/2.794/4.169 & 0.366/1.025/1.392 & 2.99 \\
  8192 & 32 & 1.423/2.846/4.260 & 0.439/1.095/1.531 & 2.78 \\
  8192 & 64 & 1.579/3.125/4.702 & 0.586/1.777/2.362 & 1.99 \\
  8192 & 128 & 1.929/3.753/5.675 & 0.896/3.200/4.076 & 1.39 \\
  16384 & 16 & 5.156/10.590/15.746 & 1.349/4.162/5.499 & 2.86 \\
  16384 & 32 & 5.320/10.800/16.118 & 1.652/4.230/5.894 & 2.73 \\
  16384 & 64 & 5.980/11.993/17.959 & 2.436/7.399/9.827 & 1.83 \\
  16384 & 128 & 7.358/14.567/21.921 & 3.588/12.543/16.094 & 1.36 \\
  32768 & 16 & 20.119/41.755/61.878 & 5.152/14.399/19.522 & 3.17 \\
  32768 & 32 & 20.753/42.579/63.357 & 6.471/16.446/22.879 & 2.77 \\
  32768 & 64 & 23.492/47.324/70.821 & 8.899/27.310/36.190 & 1.96 \\
  32768 & 128 & 29.294/57.928/87.264 & 14.300/49.573/63.947 & 1.36 \\
  65536 & 16 & 84.559/171.273/255.835 & 20.140/56.232/76.335 & 3.35 \\
  65536 & 32 & 87.328/178.228/265.548 & 25.486/64.593/90.051 & 2.95 \\
  65536 & 64 & 93.876/190.164/283.524 & 34.627/108.126/142.854 & 1.98 \\
  65536 & 128 & 113.121/227.312/339.959 & 54.549/198.086/252.637 & 1.35 \\
  \bottomrule
\end{longtable}
}

FlashAttention is faster for every configuration in this sweep, with the largest end-to-end speedups generally appearing at small $d$ and long sequence lengths: for example, FP32 at $T=65536,d=16$ improves from 255.835 ms to 76.335 ms, a 3.35x speedup. The speedup is smaller at $d=128$ because the per-tile matrix multiplications become more compute-heavy and the benefit from avoiding the full $T\times T$ attention materialization is less dominant in wall-clock time. BF16 is faster than FP32 for both implementations at long sequence lengths, but the Triton version still keeps the memory-friendly FlashAttention structure and remains consistently ahead.

\section{Distributed Data Parallel Training}

\subsection{Problem (distributed\_communication\_single\_node): Distributed Communication (Single Node)}

\begin{table}[H]
  \centering
  \caption{All-reduce runtime sweep for float32 tensors.}
  \begin{tabular}{rrrr}
    \toprule
    Tensor size & Processes/GPUs & Backend & Runtime ms \\
    \midrule
    1 MB & 2 & NCCL & -- \\
    10 MB & 2 & NCCL & -- \\
    100 MB & 2 & NCCL & -- \\
    1 GB & 2 & NCCL & -- \\
    1 MB & 4 & NCCL & -- \\
    1 MB & 6 & NCCL & -- \\
    \bottomrule
  \end{tabular}
\end{table}
\todoanswer{Add plots if useful and 2--3 sentences about data size and GPU-count effects.}

\subsection{Problem (naive\_ddp): Naive DDP}

The DDP implementation is in \codepath{cs336\_systems/ddp.py}, with adapters wired in \codepath{tests/adapters.py}. The DDP tests passed.

\subsection{Problem (naive\_ddp\_benchmarking): Naive DDP Benchmarking}

\todoanswer{Describe the 1 node x 2 GPU xl setup, total training-step time, and gradient communication time.}

\subsection{Problem (minimal\_ddp\_flat\_benchmarking): Minimal DDP with Flat Gradients Benchmarking}

\todoanswer{Report training-iteration time and communication time for a single flattened all-reduce, then compare against individual parameter all-reduces.}

\subsection{Problem (ddp\_overlap\_individual\_parameters): DDP with Overlapping Individual Parameters}

The overlapped DDP implementation is in \codepath{cs336\_systems/ddp.py}, with adapters wired in \codepath{tests/adapters.py}. The DDP tests passed.

\subsection{Problem (ddp\_overlap\_individual\_parameters\_benchmarking): DDP Overlapping Individual Parameters Benchmarking}

\begin{enumerate}
  \item \textbf{Runtime comparison.}
  \todoanswer{Report overlapped DDP time per training iteration and compare with individual and flat all-reduce baselines.}

  \item \textbf{Nsight comparison.}
  \todoanswer{Insert two Nsight screenshots showing non-overlapped vs. overlapped gradient communication.}
\end{enumerate}

\section{Optimizer State Sharding}

\subsection{Problem (optimizer\_state\_sharding): Optimizer State Sharding}

The optimizer state sharding implementation is in \codepath{cs336\_systems/sharded\_optimizer.py}, with the adapter wired in \codepath{tests/adapters.py}. The sharded optimizer tests passed.

\subsection{Problem (optimizer\_state\_sharding\_accounting): Optimizer State Sharding Accounting}

\begin{enumerate}
  \item \textbf{Peak memory breakdown.}
  \begin{table}[H]
    \centering
    \caption{Optimizer state sharding memory accounting, 1 node x 2 GPUs, xl model.}
    \begin{tabular}{lrrr}
      \toprule
      Setting & After init GiB & Before optimizer step GiB & After optimizer step GiB \\
      \midrule
      No sharding & -- & -- & -- \\
      Sharded optimizer & -- & -- & -- \\
      \bottomrule
    \end{tabular}
  \end{table}
  \todoanswer{Add 2--3 sentences breaking down parameters, gradients, and optimizer states.}

  \item \textbf{Training speed.}
  \todoanswer{Report iteration timings with and without optimizer state sharding.}

  \item \textbf{Comparison to ZeRO stage 1.}
  Our optimizer state sharding is similar to ZeRO stage 1 because each rank stores optimizer state for only a subset of parameters while full parameters and gradients remain replicated. The main difference is that our implementation shards at the granularity of whole parameter tensors using a simple round-robin assignment, which can be imbalanced, while ZeRO stage 1 partitions optimizer state more evenly across parameter elements. Our synchronization is also simpler: after each step, the owner rank broadcasts each updated parameter tensor to all other ranks, whereas optimized ZeRO implementations batch communication over flattened partitions to reduce overhead.
\end{enumerate}

\section{Fully-Sharded Data Parallel}

\subsection{Problem (fsdp): Fully-Sharded Data Parallel}

The FSDP implementation is in \codepath{cs336\_systems/fsdp.py}, with adapters wired in \codepath{tests/adapters.py}. The FSDP tests passed.

\subsection{Problem (fsdp\_accounting): FSDP Accounting}

\begin{enumerate}
  \item \textbf{Expected memory savings.}
  Since optimizer state is sharded in both optimizer-state sharding and FSDP, the additional memory savings from FSDP come primarily from sharding parameters and gradients. With two GPUs, FSDP should reduce the memory for sharded Linear and Embedding parameters and their gradients by about $2\times$ per rank relative to the optimizer-state-sharded setup. Small unsharded parameters such as normalization weights remain replicated, but they are negligible compared to the model's large weight matrices.

  \item \textbf{All-gather timing.}
  \todoanswer{Report xl two-GPU profile timings and insert Nsight screenshots showing whether all-gather completes before forward computation needs the weights.}
\end{enumerate}

\section{Analyzing Parallelism Strategies}

\subsection{Problem (alternate\_ring\_all\_reduce): Alternate Ring All-Reduce}

This algorithm takes $(N-1)S/W$ seconds, because each of the $N-1$ ring steps sends a full size-$S$ tensor from every device rather than a size-$S/N$ chunk.

\subsection{Problem (data\_parallel\_calcs): Data Parallel Calculations}

\begin{enumerate}
  \item \textbf{Backward FLOPs.}
  Per device, the backward pass takes
  \[
    \frac{12BD D_{\mathrm{FF}}}{N_{\mathrm{DP}}}
  \]
  FLOPs. There are six backward matmuls, each costing $2(B/N_{\mathrm{DP}})DD_{\mathrm{FF}}$: one for $dz$, two for $dx$, and three for the weight gradients.

  \item \textbf{Backward communication time.}
  The backward communication time is
  \[
    \frac{12(N_{\mathrm{DP}}-1)}{N_{\mathrm{DP}}}\frac{DD_{\mathrm{FF}}}{W}.
  \]
  The three weight gradients contain $3DD_{\mathrm{FF}}$ FP16 elements, or $6DD_{\mathrm{FF}}$ bytes, and ring all-reduce costs $2\frac{N_{\mathrm{DP}}-1}{N_{\mathrm{DP}}}\frac{S}{W}$.

  \item \textbf{Communication bottleneck threshold.}
  Communication is not the bottleneck while
  \[
    N_{\mathrm{DP}} \le 1 + \frac{BW}{C}.
  \]
  This follows by comparing compute time $\frac{12BD D_{\mathrm{FF}}}{N_{\mathrm{DP}}C}$ with communication time $\frac{12(N_{\mathrm{DP}}-1)DD_{\mathrm{FF}}}{N_{\mathrm{DP}}W}$ and canceling the shared model-size factors.
\end{enumerate}

\subsection{Problem (fsdp\_calcs): Fully Sharded Data Parallel Calculations}

\begin{enumerate}
  \item \textbf{Forward and backward FLOPs.}
  Per device, the forward pass takes
  \[
    \frac{6BD D_{\mathrm{FF}}}{N_{\mathrm{FSDP}}},
  \]
  and the backward pass takes
  \[
    \frac{12BD D_{\mathrm{FF}}}{N_{\mathrm{FSDP}}}.
  \]
  FSDP shards the batch across devices, so each rank computes on $B/N_{\mathrm{FSDP}}$ examples; forward has three matmuls and backward has six.

  \item \textbf{Forward and backward communication time.}
  The forward communication time is
  \[
    \frac{6(N_{\mathrm{FSDP}}-1)}{N_{\mathrm{FSDP}}}\frac{DD_{\mathrm{FF}}}{W},
  \]
  and the backward communication time is
  \[
    \frac{12(N_{\mathrm{FSDP}}-1)}{N_{\mathrm{FSDP}}}\frac{DD_{\mathrm{FF}}}{W}.
  \]
  The three weights contain $3DD_{\mathrm{FF}}$ FP16 elements, or $6DD_{\mathrm{FF}}$ bytes; forward all-gathers the weights once, while backward all-gathers the weights and reduce-scatters the gradients.

  \item \textbf{Bottleneck thresholds.}
  Both forward and backward remain compute-bound while
  \[
    N_{\mathrm{FSDP}} \le 1 + \frac{BW}{C}.
  \]
  In each case, comparing communication against compute cancels the shared $D D_{\mathrm{FF}}$ factors and leaves the same threshold.
\end{enumerate}

\subsection{Problem (tp\_calcs): Tensor Parallel Calculations}

\begin{enumerate}
  \item \textbf{Backward pass equations.}
  Since the forward pass ends with an all-reduce, every TP rank receives the same incoming $dy \in \mathbb{R}^{B \times D}$. Each rank then computes its local gradients:
  \[
    dz^{(i)} = dy (W_3^{(i)})^\top,
    \qquad
    dW_3^{(i)} = (z^{(i)})^\top dy,
  \]
  \[
    dx_2^{(i)} = dz^{(i)} \odot f(x_1^{(i)}),
    \qquad
    dx_1^{(i)} = dz^{(i)} \odot f'(x_1^{(i)}) \odot x_2^{(i)},
  \]
  \[
    dW_2^{(i)} = x^\top dx_2^{(i)},
    \qquad
    dW_1^{(i)} = x^\top dx_1^{(i)}.
  \]
  The local contribution to the input gradient is
  \[
    dx^{(i)} = dx_1^{(i)}(W_1^{(i)})^\top + dx_2^{(i)}(W_2^{(i)})^\top,
  \]
  and the final input gradient is
  \[
    dx = \mathrm{all\text{-}reduce}\left(\{dx^{(i)}\}_{i=0}^{N_{\mathrm{TP}}-1}\right).
  \]

  \item \textbf{Forward and backward FLOPs.}
  Per device, the forward pass takes
  \[
    \frac{6BD D_{\mathrm{FF}}}{N_{\mathrm{TP}}},
  \]
  and the backward pass takes
  \[
    \frac{12BD D_{\mathrm{FF}}}{N_{\mathrm{TP}}}.
  \]
  Each rank owns only $D_{\mathrm{FF}}/N_{\mathrm{TP}}$ of the FFN hidden dimension, so the three forward matmuls and six backward matmuls are reduced by a factor of $N_{\mathrm{TP}}$.

  \item \textbf{Forward and backward communication time.}
  The forward communication time is
  \[
    \frac{4BD(N_{\mathrm{TP}}-1)}{N_{\mathrm{TP}}W},
  \]
  and the backward communication time is the same:
  \[
    \frac{4BD(N_{\mathrm{TP}}-1)}{N_{\mathrm{TP}}W}.
  \]
  Forward all-reduces $y \in \mathbb{R}^{B \times D}$ and backward all-reduces $dx \in \mathbb{R}^{B \times D}$; each tensor has $BD$ FP16 elements, or $2BD$ bytes, and ring all-reduce costs $2\frac{N_{\mathrm{TP}}-1}{N_{\mathrm{TP}}}\frac{S}{W}$.

  \item \textbf{Bottleneck thresholds.}
  The forward pass remains compute-bound while
  \[
    N_{\mathrm{TP}} \le 1 + \frac{3D_{\mathrm{FF}}W}{2C},
  \]
  and the backward pass remains compute-bound while
  \[
    N_{\mathrm{TP}} \le 1 + \frac{3D_{\mathrm{FF}}W}{C}.
  \]
  These follow by comparing the shared TP all-reduce communication time $\frac{4BD(N_{\mathrm{TP}}-1)}{N_{\mathrm{TP}}W}$ against the forward compute time $\frac{6BD D_{\mathrm{FF}}}{N_{\mathrm{TP}}C}$ and backward compute time $\frac{12BD D_{\mathrm{FF}}}{N_{\mathrm{TP}}C}$.
\end{enumerate}

\subsection{Problem (fsdp\_tp\_calcs): 2D Parallelism Calculations}

\begin{enumerate}
  \item \textbf{Forward FLOPs.}
  Per device, the forward pass takes
  \[
    \frac{6BD D_{\mathrm{FF}}}{N_{\mathrm{TP}}N_{\mathrm{FSDP}}}
  \]
  FLOPs. FSDP shards the batch to $B/N_{\mathrm{FSDP}}$ examples per rank and TP shards the FFN hidden dimension to $D_{\mathrm{FF}}/N_{\mathrm{TP}}$ per rank, so each of the three forward matmuls costs $2\frac{B}{N_{\mathrm{FSDP}}}D\frac{D_{\mathrm{FF}}}{N_{\mathrm{TP}}}$ FLOPs.

  \item \textbf{Forward communication time with overlap.}
  When the FSDP-axis and TP-axis collectives can overlap, the forward communication time is
  \[
    \max\left(
      \frac{6(N_{\mathrm{FSDP}}-1)}{N_{\mathrm{FSDP}}}\frac{D D_{\mathrm{FF}}}{N_{\mathrm{TP}}W},
      \frac{4(N_{\mathrm{TP}}-1)}{N_{\mathrm{TP}}}\frac{BD}{N_{\mathrm{FSDP}}W}
    \right).
  \]
  The first term is the FSDP-axis all-gather of the three TP-sharded weights, and the second term is the TP-axis all-reduce of the FSDP-sharded output activation.

  \item \textbf{Optimal overlapped scaling threshold.}
  Let $N=N_{\mathrm{TP}}N_{\mathrm{FSDP}}$. The forward compute time is
  \[
    \frac{6BD D_{\mathrm{FF}}}{N_{\mathrm{TP}}N_{\mathrm{FSDP}}C}.
  \]
  Requiring the FSDP-axis communication term to be no larger than compute gives
  \[
    N_{\mathrm{FSDP}} \le 1 + \frac{BW}{C},
  \]
  while requiring the TP-axis communication term to be no larger than compute gives
  \[
    N_{\mathrm{TP}} \le 1 + \frac{3D_{\mathrm{FF}}W}{2C}.
  \]
  Therefore, with the optimal split across the two axes, the total device count can scale to
  \[
    N \le \left(1 + \frac{BW}{C}\right)\left(1 + \frac{3D_{\mathrm{FF}}W}{2C}\right)
  \]
  before the overlapped forward pass becomes communication-bottlenecked.

  \item \textbf{Optimal non-overlapped scaling threshold.}
  If the two axes cannot overlap, the communication time is the sum
  \[
    \frac{6(N_{\mathrm{FSDP}}-1)}{N_{\mathrm{FSDP}}}\frac{D D_{\mathrm{FF}}}{N_{\mathrm{TP}}W}
    +
    \frac{4(N_{\mathrm{TP}}-1)}{N_{\mathrm{TP}}}\frac{BD}{N_{\mathrm{FSDP}}W}.
  \]
  Comparing this to the forward compute time and multiplying through by $N_{\mathrm{TP}}N_{\mathrm{FSDP}}/D$ gives the constraint
  \[
    4B(N_{\mathrm{TP}}-1) + 6D_{\mathrm{FF}}(N_{\mathrm{FSDP}}-1)
    \le
    \frac{6BD_{\mathrm{FF}}W}{C}.
  \]
  Maximizing $N=N_{\mathrm{TP}}N_{\mathrm{FSDP}}$ under this linear constraint gives $4BN_{\mathrm{TP}}=6D_{\mathrm{FF}}N_{\mathrm{FSDP}}$, so
  \[
    N \le
    \frac{\left(\frac{6BD_{\mathrm{FF}}W}{C}+4B+6D_{\mathrm{FF}}\right)^2}{96BD_{\mathrm{FF}}}.
  \]
\end{enumerate}

\section{Leaderboard}

\subsection{Problem (leaderboard): Fastest Training Step}

\todoanswer{Report the best wall-clock time for one full forward/backward/AdamW training step for the 8B benchmark, plus the command/configuration used.}

\end{document}
